\documentclass[10pt,a4paper]{article}

%double si on veut une version prof et une version élève. simple si on veut une seule version.
\newcommand{\buildmode}{simple}

\InputIfFileExists{version.tex}{}{}
% Version (définie par le script)
\providecommand{\version}{eleve}

\usepackage{feuilleinterro}

%%%à modifier à chaque fois

\renewcommand{\niveau}{1G}
\renewcommand{\chapitre}{Fonctions, généralités}
\renewcommand{\numchap}{0}
\renewcommand{\theme}{Fonctions}

\begin{document}

% =========================================================
% PAGE 1 : VERSION A
% =========================================================

\interroversion{A}

\textbf{Exercice 1 -- Image par une fonction polynôme du second degré \hfill /1}

Soit \(g\) la fonction définie sur \(\mathbb{R}\) par \(g(x)=2x^2-3x-1\). Calculer \(g(-2)\).

\vspace{2cm}

\textbf{Exercice 2 -- Signe d'un produit de fonctions affines \hfill /4}

On considère les fonctions affines \(f\) et \(g\) définies sur \(\mathbb{R}\) par
\[
f(x)=2x-3 \qquad\text{et}\qquad g(x)=-x+4,
\]
ainsi que la fonction \(h\) définie sur \(\mathbb{R}\) par \(h(x)=f(x)\times g(x)\).

\begin{enumerate}
\item Donner les représentations graphiques approximatives des fonctions \(f\) et \(g\).

\vspace{2cm}

\item En déduire le tableau de signes de \(h\) sur \(\mathbb{R}\).

\vspace{4cm}

\item Décrire le tableau de signes obtenu.
\end{enumerate}

\newpage

% =========================================================
% PAGE 2 : VERSION B
% =========================================================

\interroversion{B}

\textbf{Exercice 1 -- Image par une fonction polynôme du second degré \hfill /1}

Soit \(g\) la fonction définie sur \(\mathbb{R}\) par \(g(x)=x^2-4x+2\). Calculer \(g(-3)\).

\vspace{2cm}

\textbf{Exercice 2 -- Signe d'un produit de fonctions affines \hfill /4}

On considère les fonctions affines \(f\) et \(g\) définies sur \(\mathbb{R}\) par
\[
f(x)=-3x+6 \qquad\text{et}\qquad g(x)=x+1,
\]
ainsi que la fonction \(h\) définie sur \(\mathbb{R}\) par \(h(x)=f(x)\times g(x)\).

\begin{enumerate}
\item Donner les représentations graphiques approximatives des fonctions \(f\) et \(g\).

\vspace{2cm}

\item En déduire le tableau de signes de \(h\) sur \(\mathbb{R}\).

\vspace{4cm}

\item Décrire le tableau de signes obtenu.
\end{enumerate}

\newpage

% =========================================================
% PAGE 3 : VERSION C
% =========================================================

\interroversion{C}

\textbf{Exercice 1 -- Image par une fonction polynôme du second degré \hfill /1}

Soit \(g\) la fonction définie sur \(\mathbb{R}\) par \(g(x)=x^2+5x-4\). Calculer \(g(-3)\).

\vspace{2cm}

\textbf{Exercice 2 -- Signe d'un produit de fonctions affines \hfill /4}

On considère les fonctions affines \(f\) et \(g\) définies sur \(\mathbb{R}\) par
\[
f(x)=4x-2 \qquad\text{et}\qquad g(x)=-2x-4,
\]
ainsi que la fonction \(h\) définie sur \(\mathbb{R}\) par \(h(x)=f(x)\times g(x)\).

\begin{enumerate}
\item Donner les représentations graphiques approximatives des fonctions \(f\) et \(g\).

\vspace{2cm}

\item En déduire le tableau de signes de \(h\) sur \(\mathbb{R}\).

\vspace{4cm}

\item Décrire le tableau de signes obtenu.
\end{enumerate}

\newpage

% =========================================================
% PAGE 4 : VERSION D
% =========================================================

\interroversion{D}

\textbf{Exercice 1 -- Image par une fonction polynôme du second degré \hfill /1}

Soit \(g\) la fonction définie sur \(\mathbb{R}\) par \(g(x)=3x^2-x+2\). Calculer \(g(-1)\).

\vspace{2cm}

\textbf{Exercice 2 -- Signe d'un produit de fonctions affines \hfill /4}

On considère les fonctions affines \(f\) et \(g\) définies sur \(\mathbb{R}\) par
\[
f(x)=-x+3 \qquad\text{et}\qquad g(x)=3x+6,
\]
ainsi que la fonction \(h\) définie sur \(\mathbb{R}\) par \(h(x)=f(x)\times g(x)\).

\begin{enumerate}
\item Donner les représentations graphiques approximatives des fonctions \(f\) et \(g\).

\vspace{2cm}

\item En déduire le tableau de signes de \(h\) sur \(\mathbb{R}\).

\vspace{4cm}

\item Décrire le tableau de signes obtenu.
\end{enumerate}

\end{document}
